\documentclass[10pt]{article}
\usepackage[utf8]{inputenc}
\usepackage[T1]{fontenc}
\usepackage[ngerman]{babel}
\usepackage[a4paper,margin=18mm]{geometry}
\usepackage{helvet}
\usepackage{xcolor}
\usepackage{colortbl}
\usepackage{tabularx}
\usepackage{longtable}
\usepackage{array}
\usepackage{enumitem}
\usepackage{fancyhdr}
\usepackage{titlesec}
\usepackage{graphicx}
\usepackage{needspace}
\renewcommand{\familydefault}{\sfdefault}
\definecolor{BEBlue}{HTML}{0066CC}
\definecolor{BEDark}{HTML}{0B2545}
\definecolor{BEOrange}{HTML}{FF9500}
\definecolor{BELight}{HTML}{F5F7FB}
\definecolor{BEPaper}{HTML}{FFF7ED}
\definecolor{BEText}{HTML}{2C3E50}
\definecolor{BEMuted}{HTML}{6C757D}
\pagestyle{fancy}
\fancyhf{}
\renewcommand{\headrulewidth}{0pt}
\lhead{\textcolor{BEDark}{\textbf{Business English in 4 Wochen}}}
\rhead{\textcolor{BEOrange}{\textbf{BE4W}}}
\cfoot{\textcolor{BEMuted}{\thepage}}
\setlength{\parindent}{0pt}
\setlength{\parskip}{5pt}
\emergencystretch=3em
\setlist[itemize]{topsep=3pt,itemsep=3pt}
\setlist[enumerate]{topsep=3pt,itemsep=3pt}
\arrayrulecolor{BEBlue!20}
\titleformat{\section}{\Large\bfseries\color{BEBlue}}{}{0pt}{}[\vspace{-2mm}\textcolor{BEOrange}{\rule{\linewidth}{1.3pt}}]
\titleformat{\subsection}{\large\bfseries\color{BEDark}}{}{0pt}{}
\newcommand{\doctype}[1]{\colorbox{BEPaper}{\parbox{0.96\linewidth}{\vspace{1mm}\textcolor{BEOrange}{\textbf{\MakeUppercase{#1}}}\vspace{1mm}}}\par\vspace{4mm}}
\newcommand{\btitle}[1]{{\Huge\bfseries\color{BEDark} #1}\par\vspace{2mm}\textcolor{BEOrange}{\rule{38mm}{2pt}}\par\vspace{3mm}}
\newcommand{\bsubtitle}[1]{{\large\color{BEText} #1}\par\vspace{7mm}}
\newcommand{\callout}[1]{\vspace{3mm}\fcolorbox{BEOrange!35}{BEPaper}{\parbox{0.94\linewidth}{\raggedright\vspace{2mm}\textcolor{BEText}{#1}\vspace{2mm}}}\vspace{4mm}}
\newcommand{\brandfooter}{\vfill\textcolor{BEMuted}{\small Wort Wach Verlag · shortcutenglish.de · Ergänzendes Material zum Buch}}
\newcommand{\coverblock}[3]{%
  \noindent\colorbox{BEDark}{\parbox{0.96\linewidth}{\raggedright\vspace{5mm}\textcolor{BEOrange}{\small\bfseries\MakeUppercase{#1}}\par\vspace{2mm}{\Huge\bfseries\textcolor{white}{#2}}\par\vspace{2mm}{\large\textcolor{white!82}{#3}}\vspace{5mm}}}\par\vspace{7mm}%
}
\newcommand{\infocard}[2]{%
  \vspace{2mm}\noindent\fcolorbox{BEBlue!18}{white}{\parbox{0.94\linewidth}{\raggedright\vspace{2.5mm}{\large\bfseries\textcolor{BEBlue}{#1}}\par\vspace{1mm}\textcolor{BEText}{#2}\vspace{2.5mm}}}\par\vspace{2mm}%
}
\newcommand{\accentcard}[2]{%
  \vspace{2mm}\noindent\fcolorbox{BEOrange!32}{BEPaper}{\parbox{0.94\linewidth}{\raggedright\vspace{2.5mm}{\large\bfseries\textcolor{BEDark}{#1}}\par\vspace{1mm}\textcolor{BEText}{#2}\vspace{2.5mm}}}\par\vspace{2mm}%
}
\newcommand{\phraseitem}[2]{%
  \noindent\fcolorbox{BEBlue!12}{white}{\parbox{0.94\linewidth}{\raggedright\vspace{1.8mm}\textcolor{BEMuted}{\small Deutsch}\par\textcolor{BEText}{#1}\par\vspace{1.2mm}\textcolor{BEOrange}{\small Englisch}\par{\bfseries\textcolor{BEDark}{#2}}\vspace{1.8mm}}}\par\vspace{2mm}%
}
\newcommand{\safechapter}[1]{\Needspace{10\baselineskip}\section{#1}}
\newcommand{\writingcard}[2]{%
  \vspace{2mm}\noindent\fcolorbox{BEOrange!32}{BEPaper}{\parbox{0.94\linewidth}{\raggedright\vspace{2.5mm}{\large\bfseries\textcolor{BEDark}{#1}}\par\vspace{#2}}}\par\vspace{2mm}%
}
\newcolumntype{Y}{>{\raggedright\arraybackslash}X}

\begin{document}
\coverblock{Bonus · Arbeitsmaterial}{Meeting Prep \& Difficult Moments}{Druckbare Arbeitsblätter für Situationen, die im normalen Training schnell zu kurz kommen}
\callout{Dieses Bonus-Workbook wiederholt nicht die Buchtheorie. Es ist für den Moment vor oder nach einem echten Meeting gedacht: vorbereiten, abdecken, laut formulieren, markieren. Shadowing und interaktives Training bleiben im Online-Coach.}

\safechapter{So nutzt du dieses Bonus-Workbook}
\accentcard{Rollenverteilung}{Das Buch erklärt die Methode. Der Coach trainiert Abruf und Wiederholung. Dieses PDF hilft dir offline bei Vorbereitung, schwierigen Gesprächsmomenten und Varianten im Ton.}
\infocard{1 · Vor dem Meeting}{Wähle drei Phrasen aus, die du wirklich verwenden willst. Schreibe sie auf die Prep-Seite und sprich sie einmal laut.}
\infocard{2 · Im Meeting}{Nutze Satzanfänge als Startbahn. Du musst nicht perfekt sein; du brauchst einen klaren Einstieg.}
\infocard{3 · Nach dem Meeting}{Markiere, welche Phrase funktioniert hat. Übertrage schwierige Sätze danach in den Coach oder in deine eigene CSV-Liste.}

\newpage
\safechapter{Quick Fix Cards: zu deutsch gedacht}
\accentcard{Ziel}{Diese Karten sind keine Grammatiklektion. Sie helfen dir, typische deutsche Direktheit und falsche Freunde in brauchbare Business-Sätze umzubauen.}
\phraseitem{Ich prüfe die Zahlen bis morgen.}{I'll check the figures by tomorrow.}
\phraseitem{Wir brauchen einen Termin mit dem Kunden.}{We need a meeting with the client.}
\phraseitem{Das aktuelle Angebot ist noch nicht final.}{The current offer is not final yet.}
\phraseitem{Ich schicke das Protokoll nach dem Meeting.}{I'll send the minutes after the meeting.}
\phraseitem{Wir brauchen einen Projektor für die Präsentation.}{We need a projector for the presentation.}
\phraseitem{Ich habe Ihre Nachricht bekommen.}{I received your message.}
\phraseitem{Ich möchte kurz den aktuellen Stand prüfen.}{I'd like to quickly check the current status.}
\phraseitem{Unsere Branche verändert sich gerade stark.}{Our industry is changing quickly at the moment.}

\newpage
\safechapter{Zu direkt → professioneller}
\accentcard{Internationaler Ton}{Die folgenden Formulierungen sind bewusst neutral gehalten. Sie funktionieren in internationalem Business English. Britisches Englisch nutzt oft mehr indirekte Höflichkeit; amerikanisches Englisch ist häufig direkter.}
\phraseitem{Schick mir den Bericht.}{Could you please send me the report when you have a moment?}
\phraseitem{Das ist falsch.}{I think we may need to adjust this part.}
\phraseitem{Du hast den Anhang vergessen.}{Could you also send the attachment?}
\phraseitem{Ich brauche sofort eine Antwort.}{Could you get back to me as soon as possible?}
\phraseitem{Das ist nicht unser Problem.}{This may be outside our current scope.}
\phraseitem{Das ergibt keinen Sinn.}{I'm not sure this part is fully clear yet.}
\phraseitem{Wir müssen die Deadline ändern.}{Would it be possible to revisit the deadline?}
\phraseitem{Du musst das heute erledigen.}{Could we make sure this is completed today?}

\safechapter{Difficult Moments Phrase Bank}
\accentcard{Warum dieser Teil exklusiv ist}{Diese Phrasen sind für heikle Momente gedacht: widersprechen, Verantwortung klären, Scope begrenzen, Verzögerungen erklären und Entscheidungen einfordern. Das ist kein normales Core-Training, sondern Meeting-Erste-Hilfe.}

\subsection{Höflich widersprechen}
\phraseitem{Ich verstehe den Punkt, aber ich sehe ein Risiko beim Zeitplan.}{I see the point, but I'm concerned about the timeline.}
\phraseitem{Ich würde das etwas anders einschätzen.}{I would look at this slightly differently.}
\phraseitem{Ich bin nicht sicher, ob das die beste Option ist.}{I'm not sure this is the best option.}
\phraseitem{Das könnte in der Praxis schwierig werden.}{That might be difficult in practice.}

\newpage
\subsection{Scope klären}
\phraseitem{Könnten wir den dringenden Teil vom langfristigen Thema trennen?}{Could we separate the urgent part from the longer-term issue?}
\phraseitem{Was ist hier im aktuellen Umfang enthalten?}{What is included in the current scope?}
\phraseitem{Das klingt nach einem zusätzlichen Arbeitspaket.}{That sounds like an additional work package.}
\phraseitem{Lassen Sie uns zuerst klären, was bis Freitag realistisch ist.}{Let's first clarify what is realistic by Friday.}

\newpage
\subsection{Verzögerungen erklären}
\phraseitem{Wir brauchen etwas mehr Zeit, um das sauber zu prüfen.}{We need a bit more time to review this properly.}
\phraseitem{Der Punkt dauert länger als ursprünglich erwartet.}{This item is taking longer than initially expected.}
\phraseitem{Ich möchte ungern eine schnelle, aber unsaubere Antwort geben.}{I'd rather not give a quick but incomplete answer.}
\phraseitem{Wir können heute einen Zwischenstand geben und morgen finalisieren.}{We can share an interim update today and finalize it tomorrow.}

\subsection{Entscheidungen einfordern}
\phraseitem{Welche Entscheidung brauchen wir heute?}{What decision do we need today?}
\phraseitem{Wer kann diesen Punkt final freigeben?}{Who can give final approval on this point?}
\phraseitem{Bis wann brauchen wir eine verbindliche Rückmeldung?}{By when do we need a firm response?}
\phraseitem{Können wir diesen Punkt jetzt entscheiden oder brauchen wir weitere Informationen?}{Can we decide this now, or do we need more information?}

\subsection{Verantwortung klären}
\phraseitem{Wer übernimmt den nächsten Schritt?}{Who will take the next step?}
\phraseitem{Ich kann den ersten Entwurf übernehmen.}{I can take care of the first draft.}
\phraseitem{Könnten Sie die technische Prüfung übernehmen?}{Could you take care of the technical review?}
\phraseitem{Lassen Sie uns die Zuständigkeiten kurz festhalten.}{Let's briefly capture the responsibilities.}

\newpage
\safechapter{UK vs US: Phrase Choices}
\accentcard{Nicht auswendig lernen}{Nutze diese Seite, um Tonunterschiede zu spüren. In internationalen Teams ist die neutrale Variante meistens die sicherste Wahl.}

\subsection{Rückfrage stellen}
\phraseitem{Neutral: Könnten Sie diesen Punkt kurz erläutern?}{Could you briefly clarify this point?}
\phraseitem{Eher britisch: Ich wollte fragen, ob Sie diesen Punkt kurz erläutern könnten.}{I was wondering if you could briefly clarify this point.}
\phraseitem{Eher amerikanisch: Können Sie diesen Punkt kurz erläutern?}{Can you briefly clarify this point?}

\subsection{Vorschlag machen}
\phraseitem{Neutral: Wir könnten den Zeitplan noch einmal prüfen.}{We could review the timeline again.}
\phraseitem{Eher britisch: Vielleicht könnten wir den Zeitplan noch einmal prüfen.}{Perhaps we could review the timeline again.}
\phraseitem{Eher amerikanisch: Lassen Sie uns den Zeitplan noch einmal prüfen.}{Let's review the timeline again.}

\subsection{Nachfassen}
\phraseitem{Neutral: Ich wollte kurz wegen der offenen Punkte nachfassen.}{I wanted to follow up on the open points.}
\phraseitem{Eher britisch: Ich wollte nur kurz wegen der offenen Punkte nachfragen.}{I just wanted to follow up on the open points.}
\phraseitem{Eher amerikanisch: Ich melde mich wegen der offenen Punkte.}{I'm following up on the open points.}

\subsection{Absage oder Einschränkung}
\phraseitem{Neutral: Das passt leider nicht in den aktuellen Zeitplan.}{Unfortunately, this does not fit into the current timeline.}
\phraseitem{Eher britisch: Ich fürchte, das könnte im aktuellen Zeitplan schwierig werden.}{I'm afraid this might be difficult within the current timeline.}
\phraseitem{Eher amerikanisch: Das ist im aktuellen Zeitplan nicht machbar.}{This is not feasible within the current timeline.}

\newpage
\safechapter{Pre-Meeting Sheet}
\accentcard{Vorbereitung auf einer Seite}{Fülle diese Seite vor einem echten Meeting aus. Ziel ist nicht Perfektion, sondern drei abrufbare Sätze, die du wirklich verwenden kannst.}
\writingcard{Mein Ziel im Meeting}{13mm}
\writingcard{Drei Phrasen, die ich aktiv verwenden will}{26mm}
\writingcard{Eine Rückfrage, die ich stellen kann}{13mm}
\writingcard{Ein höflicher Einwand, falls nötig}{13mm}
\writingcard{Follow-up-Satz für danach}{13mm}
\brandfooter
\end{document}
